%%%%%%%%%%%%%%%%%%
%% lowD-codes.tex
%%%%%%%%%%%%%%%%%%

One of the goals of the formalism in the previous chapter
is to classify all code Hamiltonians that are translationally invariant.
A satisfactory classification should be 
comprised of a complete set of invariants,
and a table of inequivalent code Hamiltonians or a machinery to obtain them.
In this chapter we present preliminary results towards this goal.

In one dimension, the classification problem is completely solved.
A particularly nice property of one dimension is
that its translation group algebra $\FF_2[x,x^{-1}]$ is a Euclidean domain,
over which any matrix is diagonal 
up to multiplications on the left and right by invertible matrices.
We show that this diagonalization is possible
even if the left multiplications are restricted to symplectic transformations.
Exploiting the finiteness of the ground field $\FF_2$,
we conclude that there are only Ising models~\cite{Yoshida2011Classification}.
An almost identical treatment appears in quantum convolutional
codes~\cite{OllivierTillich2004Convolutional,GrasslRotteler2006Convolutional}.
(A more general statement is proved by Beigi~\cite{Beigi2012OneDimension}.)

In two dimensions, we find that 
any excitation of an exact code Hamiltonian 
is described by a collection of point-like ones,
each of which is attached to a string operator.
Furthermore, it seems that one can prove a stronger statement that there are only toric codes%
~\cite{Yoshida2011Classification,BombinDuclosCianciPoulin2012}, 
although we do not complete the proof.
In three dimensions, we show that there must be point-like excitations for exact code Hamiltonians.
This implies that it is unavoidable in the translation-invariant case
to have a topological charge penetrating through the system,
for which the energy penalty is only logarithmic in the displacement.
If we assume that the degeneracy under periodic boundary conditions should be constant independent of system size,
then the energy penalty is upper bounded by a constant,
reproducing the result of Yoshida~\cite{Yoshida2011feasibility}.

% A theme underlying the proofs in this chapter is
% a geometric interpretation of coarse-graining,
% by which one changes the base ring, the translation group algebra $R$, with its subring $R'$.
% Since the coarse-grained translation group is a subgroup of the original translation group of finite index,
% $R$ is a free $R'$-module.
% Thus, the coarse-graining is an exact functor.
% Then, the algebraic set defined by the annihilator of a module 


Note that all the lemmas and theorems in this chapter are valid over qudits with prime dimensions.
We conclude this chapter with several examples.

\section{One dimension}

The group algebra $R = \FF_2[x,\bar x]$ for the one-dimensional lattice $\mathbb{Z}$
is a Euclidean domain
where the degree of a polynomial is defined to be the maximum exponent minus the minimum exponent.
(In particular, any monomial has degree $0$.)
Given two polynomials $f,g$ in $R$, one can find their $\gcd$ by the Euclid's algorithm.
It can be viewed as a column operation on the $1 \times 2$ matrix $\begin{pmatrix} f & g \end{pmatrix}$.
Similarly, one can find $\gcd$ of $n$ polynomials by column operations on $1 \times n$ matrix
\[
 \begin{pmatrix}
  f_1 & f_2 & \cdots & f_n
 \end{pmatrix} .
\]
The resulting matrix after the Euclid's algorithm will be
\[
\begin{pmatrix}
\gcd(f_1,\ldots,f_n) & 0 & \cdots & 0
\end{pmatrix} .
\]

Given a matrix $\mathbf M$ of univariate polynomials, we can apply Euclid's algorithm
to the first row and first column by elementary row and column operations
in such a way that the degree of $(1,1)$-entry $\mathbf M_{11}$ decreases unless all other
entries in the first row and column are divisible by $\mathbf M_{11}$.
Since the degree cannot decrease forever, this process must end
with all entries in the first row and column being zero except $\mathbf M_{11}$.
By induction on the number of rows or columns,
we conclude that $\mathbf M$ can be transformed to a diagonal matrix
by the elementary row and column operations.
This is known as the Smith's algorithm.

The following is a consequence of the finiteness of the ground field.
\begin{lem}
Let $\FF$ be a finite field and $S = \FF[x]$ be a polynomial ring.
Let $\phi : S \xrightarrow{f(x) \times} S$ be a $1 \times 1$ matrix such that $f(0) \ne 0$.
$\phi$ can be viewed as an $n \times n$ matrix acting on the free $S'$-module $S$
where $S' = \FF[x']$ and $x' =x^n$.
Then, for some $n \ge 1$, the matrix $\phi$ is transformed by elementary row and column operations
into a diagonal matrix with entries $1$ or $x'-1$.
The number of $x'-1$ entries in the transformed $\phi$ is equal to the degree of $f$.
\label{lem:coarse-graining-single-polynomial-in-1D}
\end{lem}
\begin{proof}
The splitting field $\tilde \FF$ of $f(x)$ is a finite extension of $\FF$.
Since $\tilde \FF$ is finite, every root of $f(x)$ is a root of $x^{n'} -1$ for some $n' \ge 1$.
Choose an integer $p \ge 1$ such that $2^p$ is greater than any multiplicity of the roots of $f(x)$.
Then, clearly $f(x)$ divides $(x^{n'}-1)^{2^p} = x^{2^p n'} - 1$.
Let $n$ be the smallest positive integer such that $f(x)$ divides $x^n-1$.%
\footnote{This part is well known, at least in the linear cyclic coding theory~\cite{MacWilliamsSloane1977}.}

Consider the coarse-graining by $S' = \FF[x']$ where $x'=x^n$.
$S$ is a free $S'$-module of rank $n$, and $(f)$ is now an endomorphism of the module $S$
represented as an $n \times n$ matrix.
Since $f(x)g(x) = x^n -1$ for some $g(x) \in \FF[x]$,
we have
\[
 A B = (x'-1) \id_n
\]
where $x' = x^n$, and $A, B$ are the matrix representation of $f(x)$ and $g(x)$, respectively, as endomorphisms.
$A$ and $B$ have polynomial entries in variable $x'$.
The determinants of $A,B$ are nonzero for their product is $(x'-1)^n \ne 0$.
Let $E_1$ and $E_2$ be the products of elementary matrices such that $A' = E_1 A E_2$ is diagonal.
Such matrices exist by the Smith's algorithm.
Put $B' = E_2^{-1} B E_1^{-1}$.
Then,
\[
 A' B' = E_1 A E_2 E_2^{-1} B E_1^{-1} = E_1 A B E_1^{-1} = (x'-1)\id_n .
\]
Since $A'$ and $I_n$ are diagonal of non-vanishing entries, $B'$ must be diagonal, too.
It follows that the diagonal entries of $A'$ divides $(x'-1)$; that is, they are $1$ or $x'-1$.

The number of entries $x'-1$ can be counted by considering $S/(f(x))$ as an $\FF$-vector space.
It is clear that $\dim_{\FF} S/(f(x)) = \deg f(x)$.
$S/(f(x)) = \coker \phi$ viewed as a $S'$-module is isomorphic to $S'^n / \im A'$,
the vector space dimension of which is precisely the number of $x'-1$ entries in $A'$.
\end{proof}

For example, consider $f(x) = x^2 + x + 1 \in S = \FF_2[x]$.
It is the primitive polynomial of the field $\FF_4$ of four elements over $\FF_2$.
Any element in $\FF_4$ is a solution of $x^4-x=0$.
Since $f(0) = 1$, we see that $n=3$ is the smallest integer such that $f(x)$ divides $x^n-1$.
As a module over $S' = \FF_2[x^3]$, the original ring $S$ is free with (ordered) basis $\{1, x, x^2\}$.
The multiplication by $x$ on $S$ viewed as an endomorphism has a matrix representation
\[
x = 
 \begin{pmatrix}
  0 & 0 & x^3 \\
  1 & 0 & 0   \\
  0 & 1 & 0   
 \end{pmatrix} .
\]
Thus, $f(x)$ as an endomorphism of $S'$-module $S$ has a matrix representation as follows.
\[
 f(x) =
 \begin{pmatrix}
  1 & x^3 & x^3 \\
  1 & 1   & x^3 \\
  1 & 1   & 1  
 \end{pmatrix} 
\cong
 \begin{pmatrix}
  1 & 0      & 0 \\
  0 & x^3 +1 & 0 \\
  0 & 0      & x^3 +1
 \end{pmatrix}
\]
Here, the second matrix is obtained by row and column operations.
There are 2 diagonal entries $x^3+1$ as $f(x)$ is of degree 2.

\begin{theorem}
If $\Lambda = \mathbb{Z}$,
any system governed by a code Hamiltonian
is equivalent to finitely many copies of Ising models,
plus some non-interacting qubits.
In particular, the topological order condition is never satisfied.
\label{thm:1D-classification}
\end{theorem}
\noindent
Yoshida~\cite{Yoshida2011Classification} arrived at a similar conclusion
assuming that the ground space degeneracy when the Hamiltonian is defined on a ring 
should be independent of the length of the ring.
If translation group is trivial,
the proof below reduces to a well-known fact that 
the Clifford group is generated by controlled-NOT, Hadamard, and Phase
gates~\cite[Proposition~15.7]{KitaevShenVyalyi2002CQC}.
The proof in fact implies that the group of all symplectic transformations
in one dimension is generated by elementary symplectic transformations
of Section~\ref{sec:symplectic-transformations}.

We will make use of the elementary symplectic transformations
and coarse-graining to deform $\sigma$ to a familiar form.
Recall that for any elementary row-addition $E$ on the upper block of $\sigma$
there is a unique symplectic transformation that restricts to $E$.

\begin{proof}
Applying Smith's algorithm to the first row and the first column of $2q \times t$ matrix $\sigma$,
one gets
\[
\begin{pmatrix}
  f_1 & 0 \\
  0   & A \\
  \hline
  g_1 & g_2 \\
  \vdots   & B  
\end{pmatrix}
\]
by elementary symplectic transformations.
Let $1 \le i < j \le q$ be integers.
If some $(1,q+j)$-entry is not divisible by $f_1$,
apply Hadamard on $j^\text{th}$ qubit to bring $(q+j)^\text{th}$ row to the upper block,
and then run Euclid's algorithm again to reduce the degree of $(1,1)$-entry.
The degree is a positive integer, so this process must end after a finite number of iteration.
Now every $(q+j,1)$-entry is divisible by $f_1$ and hence can be made to be $0$
by the controlled-NOT-Hadamard:
\[
\begin{pmatrix}
  f_1 & 0 \\
  0   & A \\
  \hline
  g_1 & g_2 \\
  0   & B  
\end{pmatrix} .
\]
Further we may assume $\deg f_1 \le \deg g_1$.
Since $\sigma^\dagger \lambda_q \sigma = 0$,
we have a commutativity condition
\[
 \bar f_1 g_1 - \bar g_1 f_1 = 0 .
\]
Write $f_1 = \alpha x^a + \cdots + \beta x^b $
and $ g_1 = \gamma x^c + \cdots + \delta x^d $
where $ a \le b$ and $c \le d$ and $\alpha,\beta,\gamma,\delta \neq 0$.
Then, $\bar f_1 g_1 = \beta \gamma x^{c-b} + \cdots + \alpha \delta x^{d-a}$.
Since $f_1 \bar g_1 = \bar f_1 g_1$, it must hold that $-(c-b)=d-a$
and $\alpha \delta = \beta \gamma$.
Since $\deg f_1 \le \deg g_1$, we have $d-b = -(c-a) \ge 0$.
The controlled-Phase $E_{1+q,1}( -(x^{d-b} + x^{c-a})\delta / \beta )$
will decrease the degree of $g_1$ by two,
which eventually becomes smaller than $\deg f_1$.
One may then apply Hadamard to swap $f_1$ and $g_1$.
Since the degree of $(1,1)$-entry cannot decrease forever,
the process must end with $g_1 = 0$.

The commutativity condition between $i^\text{th}$ ($i>1$) column and the first
is $ f_1 \bar g_i = 0 $. Since $f_1 \ne 0$, we get $g_i = 0$:
\[
\begin{pmatrix}
  f_1 & 0 \\
  0   & A \\
  \hline
  0 & 0 \\
  0 & B  
\end{pmatrix} .
\]
Continuing, we transform $\sigma$ into a diagonal matrix.
(We have shown that $\sigma$ can be transformed 
via elementary symplectic transformations to the Smith normal form.)

Now the Hamiltonian is a sum of non-interacting purely classical spin chains
plus some non-interacting qubits ($f_i = 0$).
It remains to classify classical spin chains
whose stabilizer module is generated by
\[
\begin{pmatrix}
f
\end{pmatrix}
\]
where we omitted the lower half block.
We can always choose $f=f(x)$ such that $f(x)$ has only nonnegative exponents and
$f(0)\ne 0$ since $x$ is a unit in $R$.
Lemma~\ref{lem:coarse-graining-single-polynomial-in-1D} says that
$(f)$ becomes a diagonal matrix of entries $1$ or $x'-1$
after a suitable coarse-graining followed by a symplectic transformation
and column operations.
$1$ describes the ancilla qubits,
and $x'-1 = x'+1$ does the Ising model.
\end{proof}
Note that an almost identical treatment appears in \cite{GrasslRotteler2006Convolutional}.

\section{Two dimensions}
If $D=2$, the lattice is $\Lambda = \mathbb{Z}^2$, and our base ring is $R = \FF_2 [x,\bar x, y, \bar y]$.
%%%%%%%%%%%%%%%%%%%%%

The following asserts that the local relations --- 
a few terms in the Hamiltonian that multiply to identity in a nontrivial way
as in 2D Ising model, or the kernel of $\sigma$ ---
among the terms in a code Hamiltonian,
can be completely removed for exact Hamiltonians in two dimensions~\cite{Bombin2011Structure}.
We prove a more general version.
\begin{lem}
If $G \xrightarrow{\sigma} P \xrightarrow{\epsilon} E$ is exact
over $R = F_2[ x_1, \bar x_1, \ldots, x_D, \bar x_D ]$,
There exists $\sigma' : G' \to P$ such that $\im \sigma' = \im \sigma$ and
\[
 0 \to G_{D-2} \to \cdots \to G_1 \to G' \xrightarrow{\sigma'} P \xrightarrow{\epsilon} E
\]
is an exact sequence of free $R$-modules. If $D=2$, one can choose $\sigma'$ to be injective.
\label{lem:coker-epsilon-resolution-length-D}
\end{lem}
%%%%%%%%%%%%%%%%%%%%%%%%%%%%%%%%%%%%%%%%%%%%%%%%%%%%%%%%%%%%%%%%%%%
\noindent
The lemma is almost the same as the Hilbert syzygy theorem~\cite[Corollary~15.11]{Eisenbud}
applied to $\coker \epsilon$,
which states that any finitely generated module over a polynomial ring with $n$ variables
has a finite free resolution of length $\le n$, by finitely generated free modules.
A difference is that our two maps on the far right in the resolution 
has to be related as $\epsilon = \sigma^\dagger \lambda$.
To this end, we make use of a constructive version of Hilbert syzygy theorem via Gr\"obner basis.
%%%%%%%%%%%%%%%%%%%%%%%%%%%%%%%%%%%%%%%%%%%%%%%%%%
\begin{prop}\cite[Theorem~15.10, Corollary~15.11]{Eisenbud}
Let $\{ g_1, \ldots, g_n \}$ be a Gr\"obner basis of a submodule of
a free module $M_0$ over a polynomial ring.
Then, the S-polynomials $\tau_{ij}$ of $\{ g_i \}$
in the free module $M_1 = \bigoplus_{i=1}^n S e_i$
generate the syzygies for $\{g_i\}$.
If the variable $x_1,\ldots, x_s$ are absent from the initial terms of $g_i$,
one can define a monomial order on $M_1$
such that $x_1,\ldots,x_{s+1}$ is absent from
the initial terms of $\tau_{ij}$.
If all variables are absent from the initial terms of $g_i$,
then $M_0/(g_1,\ldots,g_n)$ is free.
\label{prop:constructive-syzygy}
\end{prop}
%%%%%%%%%%%%%%%%%%%%%%%%%%%%%%%%%%%%%%%%%%%%%%%%%%%%%%%%%%%%%%%%%%
\begin{proof}[Proof of \ref{lem:coker-epsilon-resolution-length-D}]
Without loss of generality
assume that the $t \times 2q$ matrix $\epsilon$ have entries with nonnegative exponents,
so $\epsilon$ has entries in $S=\FF_2[x_1,\ldots,x_D]$.
Below, every module is over the polynomial ring $S$ unless otherwise noted.
Let $E_+$ be the free $S$-module of rank equal to $\mathrm{rank}_R~ E$.

If $g_1,\cdots, g_{2q}$ are the columns of $\epsilon$,
apply Buchberger's algorithm to obtain a Gr\"obner basis
$g_1,\cdots, g_{2q}, \ldots, g_n$ of $\im \epsilon$.
Let $\epsilon'$ be the matrix whose columns are $g_1,\ldots, g_n$.
We regard $\epsilon'$ as a map $M_0 \to E_+$.
By Proposition~\ref{prop:constructive-syzygy},
the initial terms of the syzygy generators (S-polynomials) 
$\tau_{ij}$ for $\{ g_i \}$ lacks the variable $x_1$.
Writing each $\tau_{ij}$ in a column of a matrix $\tau_1$, we have a map $\tau_1 : M_1 \to M_0$.

By induction on $D$, we have an exact sequence
\[
 M_{D} \xrightarrow{\tau_D} M_{D-1} \xrightarrow{\tau_{D-1}}
 \cdots \xrightarrow{\tau_1} M_0 \xrightarrow{\epsilon'} E_+ 
\]
of free $S$-modules, where the initial terms of columns of $\tau_D$ lack all the variables.
By Proposition~\ref{prop:constructive-syzygy} again, $M'_{D-1} = M_{D-1} / \im \tau_{\tau_D}$ is free.
Since $\ker \tau_{D-1} = \im \tau_D$, we have
\[
 0 \to M'_{D-1} \xrightarrow{\tilde \tau_{D-1}}
 \cdots \xrightarrow{\tau_1} M_0 \xrightarrow{\epsilon'} E_+ 
\]
Since $g_{2q+1},\ldots,g_n$ are $S$-linear combinations of $g_1,\ldots, g_{2q}$,
there is a basis change of $M_0$ so that the matrix representation of $\epsilon'$ becomes
\[
 \epsilon' \cong
\begin{pmatrix}
 \epsilon & 0
\end{pmatrix} .
\]
With respect to this basis of $M_0$,
the matrix of $\tau_1$ is
\[
 \tau_1 \cong
\begin{pmatrix}
 \tau_{1u} \\
 \tau_{1d}
\end{pmatrix}
\]
where $\tau_{1u}$ is the upper $2q \times t'$ submatrix.
Since $\ker \epsilon' = \im \tau_1$, 
The first row $r$ of $\tau_{1d}$ should generate $1 \in S$.
(This property is called unimodularity.)
Quillen-Suslin theorem~\cite[Chapter~XXI Theorem~3.5]{Lang}
states that there exists a basis change of $M_1$
such that $r$ becomes $\begin{pmatrix}1 & 0 & \cdots & 0 \end{pmatrix}$.
Then, by some basis change of $M_0$, one can make 
\[
\epsilon' \cong
\begin{pmatrix}
 \epsilon & 0
\end{pmatrix} ,
\quad
 \tau_{1d} \cong
\begin{pmatrix}
 1 & 0 \\
 0 & \tau'_{1d}
\end{pmatrix} .
\]
where $\tau'_{1d}$ is a submatrix.
By induction on the number of rows in $\tau_{1d}$,
we deduce that the matrix of $\tau_1$ can be brought to
\[
\epsilon' \cong
\begin{pmatrix}
 \epsilon & 0
\end{pmatrix} ,
\quad
 \tau_1 \cong
\begin{pmatrix}
 \sigma'' & \sigma' \\
 I & 0
\end{pmatrix}
\]
Note that $\epsilon \sigma'' = 0$ and $\epsilon \sigma' = 0$.
The basis change of $M_0$ by $\begin{pmatrix} I & -\sigma'' \\ 0 & I \end{pmatrix}$ gives
\[
\epsilon' \cong
\begin{pmatrix}
 \epsilon & 0
\end{pmatrix} ,
\quad
 \tau_1 \cong
\begin{pmatrix}
 0 & \sigma' \\
 I & 0
\end{pmatrix} .
\]
The kernel of $\begin{pmatrix} \sigma' \\ 0 \end{pmatrix}$ determines $\ker \tau_1 = \im \tau_2$.
Let $M'_1$ denote the projection of $M_1$ such that the sequence
\[
 0 \to M'_{D-1} \xrightarrow{\tilde \tau_{D-1}}
 \cdots \to M_2 \to M'_1 \xrightarrow{\sigma'} M'_0 \xrightarrow{\epsilon} E_+ 
\]
of free $S$-modules is exact.

Taking the ring of fractions with respect to the multiplicatively closed set 
\[
U = \{x_1^{i_1} \cdots x_D^{i_D} | i_1,\ldots,i_D \ge 0\},
\]
we finally obtain the desired exact sequence over $U^{-1}S = R$ 
with $P = U^{-1}M'_0$ and $E = U^{-1}E_+$.
Since $\im \sigma = \ker \epsilon$, we have $\im \sigma' = \im \sigma$.
\end{proof}


\begin{lem}
Let $R$ be a Laurent polynomial ring in $D$ variables over a finite field $\FF$,
and $N$ be a module over $R$.
Suppose $J = \ann_R N$ is a proper ideal such that $\dim R/J = 0$.
Then, there exists an integer $L \ge 1$ such that
\[
 \ann_{R'} N = (x_1^L -1, \ldots, x_D^L -1) \subseteq R'
\]
where $R' = \FF[x_1^{\pm L},\ldots,x_D^{\pm L}]$ is a subring of $R$.
\label{lem:zero-dim-ideal-over-finite-field}
\end{lem}

This is a variant of Lemma~\ref{lem:coarse-graining-single-polynomial-in-1D}.
The annihilator $J = \ann_R N$ is the set of all elements $r \in R$ such that
$r n = 0$ for any $n \in N$.
It is an ideal;
if $r_1, r_2 \in \ann_R N$, then
$r_1+r_2$ is an annihilator since $(r_1 + r_2)n = r_1 n + r_2 n = 0$,
and $a r_1 \in \ann_R N$ for any $a \in R$ since $(a r_1) n = a(r_1 n)=0$.
If $R' \subseteq R$ is a subring and $N$ is an $R$-module,
$N$ is an $R'$-module naturally.
Clearly, $J' = \ann_{R'} N$ is by definition equal to $(\ann_{R} N) \cap R'$.
Note that $J'$ is the kernel of the composite map $R' \hookrightarrow R \to R/J$.
Hence, we have an algebra homomorphism $\varphi' : R'/J' \to R/J$.
Although $R'$ is a subring, it is isomorphic to $R$ via the correspondence $x_i^L \leftrightarrow x_i$.
Therefore, we may view $\varphi'$ as a map $\varphi: R/I \to R/J$ for some ideal $I \subseteq R$.
It is a homomorphism such that $\varphi(x_i) = x_i^L$.
Considering the algebras as the set of all functions on the algebraic sets $V(I)$ and $V(J)$ defined by $I$ and $J$, respectively,
we obtain a map $\hat \varphi : V(J) \to V(I)$.
Intuitively, $\hat \varphi$ maps each point $(a_1,\ldots,a_D) \in \FF^D$ to $(a_1^L,\ldots,a_D^L) \in \FF^D$.
In a finite field, any nonzero element is a root of unity.
Since $\dim R/J = 0$, which means that $V(J)$ is a finite set, we can find a certain $L$
so $V(I)$ would consist of a single point. A formal proof is as follows.

\begin{proof}
Since $R$ is a finitely generated algebra over a field,
for any maximal ideal $\mm$ of $R$, the field $R/\mm$ is a finite extension of $\FF$
(Nullstellensatz~\cite[Theorem~4.19]{Eisenbud}).
Hence, $R/\mm$ is a finite field.
Since $x_i$ is a unit in $R$, the image $a_i \in R/\mm$ of $x_i$ is nonzero.
$a_i$ being an element of finite field, a power of $a_i$ is $1$.
Therefore, there is a positive integer $n$ such that
$\bb_n = (x_1^n-1,\ldots,x_D^n-1) \subseteq \mm$.
Since $x^n-1$ divides $x^{nn'}-1$,
we see that there exists $n \ge 1$ such that $\bb_n \subseteq \mm_1 \cap \mm_2$
for any two maximal ideals $\mm_1, \mm_2$.
One extends this by induction to any finite number of maximal ideals.

Since $\dim R/J = 0$, any prime ideal of $R/J$ is maximal and 
the Artinian ring $R/J$ has only finitely many maximal ideals.
$\mathrm{rad}~ J$ is then the intersection of the contractions (pull-backs)
of these finitely many maximal ideals.
Therefore, there is $n \ge 1$ such that
\[
 \bb_n \subseteq \mathrm{rad}~ J .
\]
Since $R$ is Noetherian, $(\mathrm{rad}~J)^{p^r} \subseteq J$ for some $r \ge 0$
where $p$ is the characteristic of $\FF$. Hence, we have 
\[
\bb_{np^r} \subseteq \bb_n^{p^r} \subseteq (\mathrm{rad}~J)^{p^r} \subseteq J.
\]

Let $L = np^r$.
If $R' = \FF[x_1^L,\bar x_1^L,\ldots,x_D^L,\bar x_D^L]$,
$\ann_{R'} N$ is nothing but $J \cap R'$.
We have just shown $\bb_L \cap R' \subseteq J \cap R'$.
Since $J$ is a proper ideal, we have $1 \notin J \cap R'$.
Thus, $\bb_L \cap R'= J \cap R'$ since $\bb_L \cap R'$ is maximal in $R'$.
\end{proof}


\begin{theorem}
For any two-dimensional degenerate exact code Hamiltonian,
there exists an equivalent Hamiltonian such that
\[
\ann \coker \epsilon = (x-1,y-1).
\]
Thus, $\coker \epsilon$ is a torsion module.
\label{thm:structure-2d-ann-coker-epsilon}
\end{theorem}
The content of Theorem~\ref{thm:structure-2d-ann-coker-epsilon}
is presented in \cite{Bombin2011Structure}.
We will comment on it after the proof.
\begin{proof}
By Lemma~\ref{lem:coker-epsilon-resolution-length-D},
we can find an equivalent Hamiltonian such that the generating map $\sigma$ for its stabilizer module
is injective:
\[
 0 \to G \xrightarrow{\sigma} P .
\]
Let $t$ be the rank of $G$.
The exactness condition says
\[
 0 \to G \xrightarrow{\sigma} P \xrightarrow{\epsilon} E
\]
is exact where $\epsilon = \sigma^\dagger \lambda_q$ and $E$ has rank $t$.
Applying Proposition~\ref{prop:exact-sequence},
since $\overline{ I(\sigma) } = I(\epsilon)$ 
and hence in particular $\codim I(\sigma) = \codim I(\epsilon)$,
we have that $q=t$ and $\codim I(\epsilon) \ge 2$ if $I(\epsilon) \ne R$.
But, $I(\epsilon) \ne R$ by Corollary~\ref{cor:unit-characteristic-ideal-means-nondegeneracy}.

Since $q=t$, $I(\epsilon)$ is equal to the initial Fitting ideal,
and therefore has the same radical as the annihilator of $\coker \epsilon = E/ \im \epsilon$.
(See \cite[Proposition~20.7]{Eisenbud} or \cite[Chapter~XIX Proposition~2.5]{Lang}.)
In particular, $\dim R / (\ann \coker \epsilon) = 0$.
Apply Lemma~\ref{lem:zero-dim-ideal-over-finite-field} to conclude the proof.
\end{proof}

An interpretation of the theorem is the following.
For systems of qubits, Theorem~\ref{thm:structure-2d-ann-coker-epsilon} says that
$x+1$ and $y+1$ are in $\ann \coker \epsilon$.
In other words, any element $v$ of $E$ is a charge,
and a pair of $v$'s of distance 1 apart
can be created by a local operator.
Equivalently, $v$ can be translated by distance 1 by the local operator.
Since translation by distance 1 generates all translations of the lattice,
we see that any excitation can be moved through the system by some sequence of local operators.
This is exactly what happens in the 2D toric code:
Any excited state is described by a configuration of magnetic and electric charge,
which can be moved to a different position by a string operator.

Moreover, since $(x-1,y-1) = \ann \coker \epsilon$,
the action of $x,y \in R$ on $\coker \epsilon$ is the same as the identity action.
Therefore, the $R$-module $\coker \epsilon$ is completely determined
up to isomorphism by its dimension $k$ as an $\FF_2$-vector space.
In particular, $\coker \epsilon$ is a finite set,
which means there are finitely many charges.
The module $K(L)$ of Pauli operators acting on the ground space (logical operators),
can be viewed as $K(L) = \Tor_1(\coker \epsilon, R/\bb_L)$.
Thus, $K(L)$ is determined by $k$ up to $R$-module isomorphisms.
This implies that the translations of a logical operator are all equivalent.
It is not too obvious at this moment whether the symplectic structure,
or the commutation relations among the logical operators,
of $K(L)$ is also completely determined.

Yoshida~\cite{Yoshida2011Classification} argued a similar result assuming
that the ground-state degeneracy should be independent of system size.
Bombin~\cite{Bombin2011Structure} later claimed without the constant degeneracy assumption
that one can choose locally independent stabilizer generators in a `translationally invariant way'
in two dimensions, for which Lemma~\ref{lem:coker-epsilon-resolution-length-D} is a generalization,
and that there are finitely many topological charges,
which is immediate from Theorem~\ref{thm:structure-2d-ann-coker-epsilon} since $\coker \epsilon$ is a finite set.
The claim is further strengthened assuming extra conditions by Bombin~\emph{et al.}~\cite{BombinDuclosCianciPoulin2012},
which can be summarized by saying
that $\sigma$ is a finite direct sum of $\sigma_\text{2D-toric}$ in Example~\ref{eg:2d-toric}.

\begin{rem}
\label{rem:long-string}
Although the strings are capable of moving charges on the lattice,
it could be very long compared to the interaction range.
Consider
\[
 \epsilon_p = 
\begin{pmatrix}
 p(x) & p(y) & 0         & 0          \\
 0    & 0    & p(\bar y) & -p(\bar x)
\end{pmatrix}
\]
where $p$ is any polynomial. It defines an exact code Hamiltonian.
For instance, the choice $p(t) = t-1$ reproduces the 2D toric code of Example~\ref{eg:2d-toric}.
Now let $p(t)$ be a primitive polynomial of the extension field $\FF_{2^w}$ over $\FF_2$.
$p(t)$ has coefficients in the base field $\FF_2$ and factorizes in $\FF_{2^w}$
as $p(t) = (t-\theta)(t-\theta^2)(t-\theta^{2^2})\cdots(t-\theta^{2^{w-1}})$.
(See \cite[Chapter~V Section~5]{Lang}.)
The multiplicative order of $\theta$ is $N = 2^w -1$.
The degree $w$ of $p(t)$ may be called the interaction range.
If the charge $e = \begin{pmatrix}1 & 0 \end{pmatrix}^T$ at $(0,0)\in \ZZ^2$
is transported to $(a,b) \in \ZZ^2 \setminus \{ (0,0) \}$ by some finitely supported operator,
we have $(x^a y^b -1)e \in \im \epsilon$.
That is, $x^a y^b -1 \in (p(x),p(y))$.
Substituting $x \mapsto \theta$ and $y \mapsto \theta^{2^m}$,
we see that $\theta^{a+2^m b} = 1$ or $a+2^m b \equiv 0 \pmod N$ for any $m \in \ZZ$.
In other words, $a \equiv -b \equiv -2b \pmod N$. It follows that $|a| + |b| \ge N$.%
\footnote{In case of qudits with prime dimensions $p$,
the lower bound will be $N/(p-1)$.}
Therefore, the length of the string segment transporting a charge is exponential in the interaction range $w$.
\end{rem}

\begin{example}[Wen plaquette~\cite{Wen2003Plaquette}]
\label{eg:wen-plaquette}
This model consists of a single type of interaction ($t=q=1$)
\[
\xymatrix@!0{
X \ar@{-}[r] & Y \ar@{-}[d] \\
Y \ar@{-}[u] & X \ar@{-}[l]
} \quad \quad
\sigma_\text{Wen} = 
\begin{pmatrix}
 1 + x + y + xy \\
\hline
 1 + xy
\end{pmatrix}
\]
where $X,Y$ are abbreviations of $\sigma_x,\sigma_y$.
It is known to be equivalent to the 2D toric code.
Take the coarse-graining given by $R' = \FF_2[x',y',\bar x', \bar y']$ where
\[
 x' = x \bar y, \quad \quad y' = y^2 .
\]
(The coarse-graining considered in this example
is intended to demonstrate a non-square blocking of the old lattice
to obtain a `tilted' new lattice, and is by no means special.)
As an $R'$-module, $R$ is free with basis $\{ 1, y \}$.
With the identification $R = (R' \cdot 1) \oplus (R' \cdot y)$,
we have $x \cdot 1 = x' \cdot y$, $x \cdot y = x'y' \cdot 1$,
and $y \cdot 1 = 1 \cdot y$, $y \cdot y = y' \cdot 1$.
Hence, $x$ and $y$ act on $R'$-modules
as the matrix-multiplications on the left:
\[
x \mapsto 
\begin{pmatrix}
0  & x'y'\\
x' &   0
\end{pmatrix},
\quad
y \mapsto
\begin{pmatrix}
0 & y'\\
1 & 0
\end{pmatrix}.
\]
Identifying
\[
 R^n = [ (R' \cdot 1) \oplus (R' \cdot y) ] \oplus \cdots \oplus [(R' \cdot 1) \oplus (R' \cdot y)] ,
\]
our new $\sigma$ on the coarse-grained lattice becomes
\[
 \sigma' =
\begin{pmatrix}
1+x'y' & y'+x'y'\\
1+x'   & 1+x'y'\\
\hline
1+x'y' & 0 \\
0      & 1+x'y'
\end{pmatrix}.
\]
By a sequence of elementary symplectic transformations, we have
\begin{align*}
\sigma'
&
\xrightarrow[E_{1,3}(1)]{E_{2,4}(1)}
%
\begin{pmatrix}
0      & y'+x'y' \\
1+x'   & 0       \\
1+x'y' & 0       \\
0      & 1+x'y'
\end{pmatrix} 
%
\xrightarrow[E_{3,2}(y')]{E_{4,1}(\bar y')}
%
\begin{pmatrix}
0      & y'+x'y' \\
1+x'   & 0       \\
1+y' & 0       \\:
0      & x'y'+x'
\end{pmatrix}
%
\\
& \xrightarrow[\times \bar x' \bar y']{\text{col.2}}
%
\begin{pmatrix}
0      & 1 + \bar x' \\
1+x'   & 0       \\
1+y' & 0       \\
0      & 1+\bar y'
\end{pmatrix}
%
\xrightarrow{1 \leftrightarrow 3}
%
\begin{pmatrix}
1+y' & 0     \\
1+x'   & 0   \\
0      & 1+ \bar x'\\
0      & 1 + \bar y'\\
\end{pmatrix},
\end{align*}
which is exactly the 2D toric code.
\hfill $\Diamond$
\end{example} % Wen Plaquette



%(include?) \input{extracting-toric-codes}


\section{Three dimensions}
\label{sec:3d}

In the previous section, we derived a consequence of the exactness of code Hamiltonians.
The two-dimensional Hamiltonian was special 
so we were able to characterize the behavior of the charges more or less completely.
Here, we prove a weaker property of three dimensions that there must exist a nontrivial charge
for any exact code Hamiltonian.
It follows from Theorems~\ref{thm:charge-equals-torsion},%
\ref{thm:logarithmic-energy-barrier-for-fractal-operator}
that such a charge can spread through the system
by surmounting the logarithmic energy barrier.
\begin{lem}
Suppose $D=3$,
\[
  0 \to G_1 \xrightarrow{\sigma_1} G \xrightarrow{\sigma} P \xrightarrow{\epsilon=\sigma^\dagger \lambda_q} E
\]
is exact, and $I(\sigma) \subseteq \mm = (x-1,y-1,z-1)$.
Then, $\coker \epsilon$ is not torsion-free.
\label{lem:char-ideal-with-support-on-origin-implies-fractal}
\end{lem}
\begin{proof}
Suppose on the contrary $\coker \epsilon$ is torsion-free.
We have an exact sequence
\[
 0 \to G_1 \xrightarrow{\sigma_1} G \xrightarrow{\sigma} P \xrightarrow{\epsilon} E \to E'.
\]
If $G_1 = 0$, Proposition~\ref{prop:injective-sigma-implies-fractal} implies the conclusion.
So we assume $G_1 \ne 0$, and therefore we have $I(\sigma_1) = R$ by Proposition~\ref{prop:exact-sequence}.

Let us localize the sequence at $\mm$, so $I(\sigma_1)_\mm = R_\mm$.
Since $\rank (G_1)_\mm = \rank (\sigma_1)_\mm$, the matrix of $(\sigma_1)_\mm$ becomes
\[
 (\sigma_1)_\mm =
\begin{pmatrix}
 0 \\
 I
\end{pmatrix}
\]
for some basis of $(G_1)_\mm$ and $G_\mm$. See the proof of 
Lemma~\ref{lem:associated-maximal-localized-homology}.
In other words, there is an invertible matrix $B \in \mathrm{GL}_{ t \times t}( R_\mm )$ such that
\[
 \sigma_\mm B = 
\begin{pmatrix}
 \tilde \sigma & 0
\end{pmatrix}
\]
where $\tilde \sigma$ is the $2q \times t'$ submatrix.
Note that the antipode map is a well-defined automorphism of $R_\mm$
since $\overline \mm = \mm$.

Since $\epsilon = \sigma^\dagger \lambda_q$, we have
\begin{equation}
 B^\dagger \epsilon_\mm = 
\begin{pmatrix}
 \tilde \sigma^\dagger \\
 0
\end{pmatrix} \lambda_q
=
\begin{pmatrix}
 \tilde \sigma^\dagger \lambda_q \\
 0
\end{pmatrix} .
\label{eq:tilde-epsilon-sigma-dagger}
\end{equation}
Therefore, we get a new exact sequence
\[
   0 \to G' \xrightarrow{\tilde \sigma} P_\mm \xrightarrow{\tilde \epsilon = \tilde \sigma^\dagger \lambda_q} R_\mm^{t'}
\]
where $G' = G_\mm / \im (\sigma_1)_\mm$ is a free $R_\mm$-module and $t' = \rank G'$.
It is clear that 
$\rank \tilde \epsilon = \rank \tilde \sigma$.
Setting $S = R_\mm$ in Proposition~\ref{prop:injective-sigma-implies-fractal} implies that
$\coker \tilde \epsilon$ is not torsion-free.
But, since we are assuming $\coker \epsilon_m$ is torsion-free,
$\coker \tilde \sigma^\dagger$ is also torsion-free by Eq.~\eqref{eq:tilde-epsilon-sigma-dagger}.
This is a contradiction.
\end{proof}



\begin{theorem}
For any three-dimensional, degenerate and exact code Hamiltonian,
there exists a fractal generator.
\label{thm:fractal-exists-in-3D}
\end{theorem}
\begin{proof}
By Lemma~\ref{lem:coker-epsilon-resolution-length-D}, there exists an equivalent Hamiltonian
such that
\[
  0 \to G_1 \xrightarrow{\sigma_1} G \xrightarrow{\sigma} P \xrightarrow{\epsilon=\sigma^\dagger \lambda_q} E
\]
is exact. The existence of a fractal generator is a property of the equivalence class
by Proposition~\ref{prop:fractal-property-of-eq-class}.
If we can find a coarse-graining such that $I(\sigma') \subseteq (x'-1,y'-1,z'-1)$,
then Lemma~\ref{lem:char-ideal-with-support-on-origin-implies-fractal} shall imply the conclusion.


Recall that $\epsilon_L$ and $\sigma_L$ denote the induced maps
by factoring out $\bb_L = (x^L-1,y^L-1,z^L-1)$. See Sec.~\ref{sec:degeneracy}.
There exists $L$ such that $K(L) = \ker \epsilon_L / \im \sigma_L \ne 0$ by 
Corollary~\ref{cor:unit-characteristic-ideal-means-nondegeneracy}.
Consider the coarse-grain by $x'=x^L,~ y'=y^L,~ z'=z^L$.
Let $R' = F_2[x'^{\pm 1}, y'^{\pm 1}, z'^{\pm 1}]$ denote the coarse-grained base ring.
If $K'(L')$ denotes $\ker \epsilon'_{L'} / \im \sigma'_{L'}$ as $R'$-module,
we see that $K'(1) = K(L)$ as $\FF_2$-vector space.
In particular, $K'(1) \ne 0$.
Put $\mm = (x'-1,y'-1,z'-1) = \bb'_1 \subseteq R'$.
Then, $K'(1)_{\mm} = K'(1) \ne 0$.
By Lemma~\ref{lem:associated-maximal-localized-homology},
we have $I(\sigma') \subseteq \mm$.
\end{proof}

Yoshida argued that when the ground-state degeneracy is constant independent of system size
there exists a string operator~\cite{Yoshida2011feasibility}.
To prove it, we need an algebraic fact.
\begin{prop}
Let $M$ be a finitely presented $R$-module, and $T$ be its torsion submodule.
Let $I$ be the first non-vanishing Fitting ideal of $M$.
Then,
\[
 \rad I \subseteq \rad \ann T .
\]
\label{prop:support-torsion-fitting-ideal}
\end{prop}
\begin{proof}
Let $\pp$ be any prime ideal of $R$ such that $I \not\subseteq \pp$.
By the calculation of the proof of Lemma~\ref{lem:associated-maximal-localized-homology},
$M_\pp$ is a free $R_\pp$-module, and hence is torsion-free.
Since $T$ is embedded in $M$, it follows that $T_\pp = 0$, or equivalently, $\ann T \not\subseteq \pp$.
Since the radical of an ideal is the intersection of all primes containing it~\cite[Proposition 1.8]{AtiyahMacDonald},
the claim is proved.
\end{proof}

\begin{cor}
Let $T$ be the set of all point-like charges modulo locally created ones
of a degenerate and exact code Hamiltonian in three dimensions
of characteristic dimension zero.
Then, one can coarse-grain the lattice such that
\[
 \ann T = (x-1,y-1,z-1) .
\]
\label{cor:annT-3d-characteristic-dimension-0}
\end{cor}
\noindent
The corollary says that any point-like charge is attached to strings and is able to move freely through the lattice.
The condition is implied by Lemma~\ref{lem:k-growing-sequence-L} 
if the ground-state degeneracy is constant independent of the system size
when defined on a periodic lattice.
%%%%%%
\begin{proof}
By Theorem~\ref{thm:charge-equals-torsion}, $T$ is the torsion submodule of $\coker \epsilon$.
By Theorem~\ref{thm:fractal-exists-in-3D}, $T$ is nonzero.
Setting $M = \coker \epsilon$ in Proposition~\ref{prop:support-torsion-fitting-ideal},
the associated ideal $I_q(\epsilon)$ is the first non-vanishing Fitting ideal of $M$.
Since $\dim R / I_q(\epsilon) = 0$ by assumption, we have $\dim R / \ann T = 0$.
Lemma~\ref{lem:zero-dim-ideal-over-finite-field} implies the claim.
\end{proof}

\begin{example}[Toric codes in higher dimensions]
\label{eg:highD-toric-codes}
Any higher-dimensional toric code can be treated similarly as for the two-dimensional case.
In three dimensions one associates each site with $q=3$ qubits.
It is easily checked that
\[
 \sigma_\text{3D-toric} = 
\begin{pmatrix}
1 + \bar x & 0   & 0      & 0    \\
1 + \bar y & 0   & 0      & 0    \\
1 + \bar z & 0   & 0      & 0    \\
\hline
0          & 0   & 1+z    & 1+y  \\
0          & 1+z & 0      & 1+x  \\
0          & 1+y & 1+x    & 0    \\
\end{pmatrix} .
\]

Both two- and three-dimensional toric codes have the property 
that $\coker \epsilon$ is not torsion-free.
However, in two dimensions any element of $E$ is a physical charge,
whereas in three dimensions $E$ contains physically irrelevant elements.
Note that in both cases, $1+x$ and $1+y$ are fractal generators.
Being consisted of two terms, they generate the `string operators.'


The 4D toric code~\cite{DennisKitaevLandahlEtAl2002Topological}
has $\sigma_x$-type interaction and $\sigma_z$-type interaction.
Originally the qubits are placed on every plaquette of 4D hypercubic lattice;
instead we place $q=6$ qubits on each site.
The generating map $\sigma$ for the stabilizer module is written as
a $12 \times 8$-matrix ($t=8$)
\[
\sigma_\text{4D-toric} =
\begin{pmatrix}
 \sigma_X & 0 \\
 0        & \sigma_Z
\end{pmatrix}
\]
where
\begin{align*}
 \sigma_X &= 
\begin{pmatrix}
1+y & 1+x & 0   & 0   \\
1+w & 0   & 0   & 1+x \\
1+z & 0   & 1+x & 0   \\
0   & 1+z & 1+y & 0   \\
0   & 1+w & 0   & 1+y \\
0   & 0   & 1+w & 1+z
\end{pmatrix}, \\
\bar \sigma_Z &=
\begin{pmatrix}
0   & 0   & 1+w & 1+z \\
0   & 1+z & 1+y & 0   \\
0   & 1+w & 0   & 1+y \\
1+w & 0   & 0   & 1+x \\
1+z & 0   & 1+x & 0   \\
1+y & 1+x & 0   & 0
\end{pmatrix}.
\end{align*}
Note the bar on $\sigma_Z$.

Theorem~\ref{thm:fractal-exists-in-3D} does not prevent
the absence of a fractal generator in four or higher dimensions.
Indeed, this 4D toric code lacks any fractal generator.
To see this, it is enough to consider $\sigma_Z$
since $\overline{\coker \sigma_X^\dagger} \cong \coker \sigma_Z^\dagger$ as $R_4$-modules,
where $R_4 = \FF_2[x^{\pm 1},y^{\pm 1},z^{\pm 1},w^{\pm 1}]$.
If
\[
\epsilon_1 =
 \begin{pmatrix}
1+x & 1+y & 1+z & 1+w
 \end{pmatrix} : R_4^4 \to R_4,
\]
then
\[
 R_4^6 \xrightarrow{\sigma_Z^\dagger} R_4^4 \xrightarrow{\epsilon_1} R_4
\]
is exact.
(A direct way to check it is to compute S-polynomials~\cite[Chapter~15]{Eisenbud}
of the entries of $\epsilon_1$,
and to verify that they all are in the rows of $\sigma_Z$.)
Hence, $\coker \sigma_Z^\dagger$ is torsion-free by Proposition~\ref{prop:nofractal-torsionless}.


For the toric codes in any dimensions, $\sigma$ has nonzero entries of form $x_i-1$.
The radical of the associated ideal $I(\sigma)$ is equal to $\mm = (x_1-1,\ldots, x_D-1)$.
So $\mm$ is the only maximal ideal of $R$ that contains $I(\sigma)$.
The characteristic dimension is zero.
If $2 \nmid L$, since $(\bb_L)_{\mm} = \mm_{\mm}$,
$(\sigma_L)_\mm$ is a zero matrix.
Any other localization of $\sigma_L$ does not contribute to $\dim_{\FF_2} K(L)$
by Lemma~\ref{lem:associated-maximal-localized-homology}.
Therefore, if $2 \nmid L$, $K(L)$ has constant vector space dimension independent of $L$.

There is a more direct way to compute the $R$-module $K(L)$.
For the three-dimensional case, consider a free resolution of $R_3/\mm$,
where $R_3 = \FF_2[x^{\pm 1},y^{\pm 1},z^{\pm 1}]$,
as
\[
 0 \to 
R_3^1 
\xrightarrow{
\partial_3 = 
\begin{pmatrix} a \\ b \\ c \end{pmatrix}
}
R_3^3
\xrightarrow{  
\partial_2 =
\begin{pmatrix}
 0 & -c & b  \\
 c & 0 & -a  \\
 -b & a & 0   
\end{pmatrix}
}
R_3^3
\xrightarrow{
\partial_1 =
\begin{pmatrix}
a & b  & c
\end{pmatrix}
}
R_3^1
\to
R_3/\mm
\to 0
\]
where $a=x-1$, $b=y-1$, and $c=z-1$.
We see that
\begin{equation}
 \sigma_\text{3D-toric} = \bar \partial_3 \oplus \partial_2 ,
 \quad \text{and} \quad
 \epsilon_\text{3D-toric} = \bar \partial_2 \oplus \partial_1  .
\label{eq:resolution-F}
\end{equation}
Therefore,
\begin{align*}
 K(L)_\text{3D-toric} \cong \Tor_1(\coker \epsilon_\text{3D-toric},R_3/\bb_L) 
&\cong \Tor_2( \overline{ R_3/\mm}, R_3/\bb_L ) \oplus \Tor_1( R_3/\mm, R_3/\bb_L ).
\end{align*}
Using $\Tor(M,N) \cong \Tor(N,M)$ and the fact that a resolution of $R_3/\bb_L$ is Eq.~\eqref{eq:resolution-F}
with $a,b,c$ replaced by $x^L-1,y^L-1,z^L-1$, respectively,
we have 
\[
 \Tor_i(R_3/\mm,R_3/\bb_L) \cong \Tor_i(R_3/\mm,R_3/\mm) \cong (\FF_2)^{_3 C _i}
\]
for each $0 \le i \le 3$.
Therefore, $K(L)_\text{3D-toric} \cong (\FF_2)^{_3 C _2} \oplus (\FF_2)^{_3 C _1} \cong (\FF_2)^6$.
The four-dimensional case is similar:
\[
 K(L)_\text{4D-toric} \cong \Tor_2 (R_4/\mm, R_4/\bb_L) \oplus \Tor_2 (\overline{R_4/\mm}, R_4/\bb_L) \cong \left( (\FF_2)^{_4 C _2} \right)^2 .
\]
The calculation here is closely related
to the cellular homology interpretation of toric codes.
\hfill $\Diamond$
\end{example} % n dimensional toric codes

\begin{example}[Chamon model~\cite{Chamon2005Quantum,BravyiLeemhuisTerhal2011Topological}]
This three-dimensional model consists of single type of term in the Hamiltonian.
The generating map is
\[
\sigma_\text{Chamon} =
\begin{pmatrix}
 x+\bar x + y + \bar y \\
\hline
 z+\bar z + y + \bar y
\end{pmatrix}.
\]
Since
\[
 \sigma^\dagger \lambda_1
\begin{pmatrix}
0 \\ 
1
\end{pmatrix}
=
(1+ x \bar y)
\begin{pmatrix}
 0 \\
\bar x + y
\end{pmatrix},
\]
$1+x \bar y$ is a fractal generator.
Consisted of two terms, it generates a string operator.
The degeneracy can be calculated using Corollary~\ref{cor:k-formulas}.
Assume all the three linear dimensions of the system are even.
Put
\[
 S = R / (x + \bar x + y + \bar y, z + \bar z + y + \bar y, x^{2l} -1, y^{2m}-1, z^{2n} -1 ).
\]
Then, the $\log_2$ of the degeneracy is $k = \dim_{\FF_2} S$.
In $S$, we have $x + \bar x = y + \bar y = z + \bar z$.
Since $S$ has characteristic 2, it holds that
\[
 w^{p+1} + w^{-p-1} = (w + w^{-1})(w^p + w^{p-2} + \cdots + w^{-p})
\]
for $p \ge 1$ and $w = x,y,z$. By induction on $p$, we see that $w^{p} + w^{-p}$
is a polynomial in $w + w^{-1}$. Therefore,
\[
 x^p + \bar x ^p = y^p + \bar y^p = z^p + \bar z^p
\]
for all $p \ge 1$ in $S$.
Put $g = \gcd(l,m,n)$.
Since $x^l + x^{-l} = y^m + y^{-m} = z^n + z^{-n} = 0$ in $S$,
we have $x^g + x^{-g} = y^g + y^{-g} = z^g + z^{-g} = 0$.

Applying Buchberger's criterion with respect to the lexicographic order in which $x \prec y \prec z$,
we see that
\[
 S = \FF_2[x,y,z] / (z^2 + zx^{2l-1}+ zx+1, y^2 + yx^{2l-1}+yx+1, x^{2g}+1)
\]
is expressed with a Gr\"obner basis. Therefore,
\[
 k = \dim_{\FF_2} S = 8 \gcd(l,m,n).
\]
\hfill $\Diamond$
\label{eg:ChamonModel}
\end{example}

\begin{example}[Levin-Wen fermion model~\cite{LevinWen2003Fermions}]
\label{eg:Levin-Wen-fermion-model}
The 3-dimensional model is originally defined 
in terms of Hermitian bosonic operators $\{ \gamma^{ab} \}_{a,b=1,\ldots,6}$,
squaring to identity if nonzero,
such that $\gamma^{ab}=-\gamma^{ba}$,
$[\gamma^{ab},\gamma^{cd}]=0$ if $a,b,c,d$ are distinct,
and
$\gamma^{ab}\gamma^{bc}=i\gamma^{ac}$ if $a \neq c$.
An irreducible representation is given by Pauli matrices acting on $\mathbb{C}^2 \otimes \mathbb{C}^2$,
and their commuting Hamiltonian fits nicely into our formalism.
The model was proposed to demonstrate that the point-like excitations may actually be fermions.
\begin{align*}
 \sigma_\text{Levin-Wen} &=
\begin{pmatrix}
 1+z & 1+z & x+y \\
 y+y z & x+x z & x+y \\
 y+z & 1+x & 1+x \\
 y+z & z+x z & y+x y
\end{pmatrix} \\
 \epsilon_\text{Levin-Wen} &=
\begin{pmatrix}
 y+z & y+z & y+y z & 1+z \\
 z+x z & 1+x & x+x z & 1+z \\
 y+x y & 1+x & x+y & x+y
\end{pmatrix}
\end{align*}
Here we multiplied the rows of $\epsilon_\text{Levin-Wen}$ by suitable monomials to avoid negative exponents.
One readily verifies that $\ker \epsilon_\text{Levin-Wen} = \im \sigma_\text{Levin-Wen}$.
The model is symmetric under the spatial rotation by $\pi/3$ about $(1,1,1)$ axis.
Indeed, if one changes the variables as $x \mapsto y \mapsto z \mapsto x$ and apply a symplectic transformation
\begin{equation}
 \omega = 
 \begin{pmatrix}
 1 & 0 & 1 & 0 \\
 0 & 1 & 0 & 1 \\
 1 & 0 & 0 & 0 \\
 0 & 1 & 0 & 0
 \end{pmatrix} :
 \begin{cases}
   XI \mapsto YI \\
   IX \mapsto IY \\
   ZI \mapsto XI \\
   IZ \mapsto IX
\end{cases} ,
\label{eq:symplectic-Levin-Wen}
\end{equation}
then $\sigma_\text{Levin-Wen}$ remains the same up to permutations of columns.


The torsion submodule $T$ of $C = \coker \epsilon_\text{Levin-Wen}$,
which describes the point-like charges according to Theorem~\ref{thm:charge-equals-torsion},
is
\begin{equation}
 T = R \cdot \begin{pmatrix} 1+y \\ 1+x \\ 0 \end{pmatrix} .
\label{eq:Levin-Wen-torsion}
\end{equation}
In order to see this, first shift the variables $a = x+1, b=y+1, c=z+1$.
Then, $\epsilon_\text{Levin-Wen}$ becomes
\[
\epsilon_\text{Levin-Wen} =
\begin{pmatrix}
 b+c & b+c & c+b c & c \\
 a+a c & a & c+a c & c \\
 a+a b & a & a+b & a+b
\end{pmatrix}
 =: \phi
\]
We will verify that $N = C / T$ is torsion-free.
A presentation of $N = \coker \phi'$ is obtained by joining the generator of $T$ to the matrix $\phi$.
\[
 \phi' = 
\begin{pmatrix}
 b+c & b+c & c+b c & c & b \\
 a+a c & a & c+a c & c & a \\
 a+a b & a & a+b & a+b & 0
\end{pmatrix}
\]
Column operations of $\phi'$ give
\[
 \phi' \cong
\begin{pmatrix}
 0 & c & b & 0 & 0 \\
 c & 0 & a & 0 & 0 \\
 b & a & 0 & 0 & 0
\end{pmatrix} = 
\begin{pmatrix}
 \partial_2 & 0 & 0
\end{pmatrix}
\]
where $\partial_2$ is from Eq.~\eqref{eq:resolution-F}.
Therefore, $\phi'$ generates the kernel of $\partial_1$,
and by Proposition~\ref{prop:nofractal-torsionless}, $N = \coker \phi'= \coker \partial_2$ is torsion-free.

The torsion submodule $T$ of $C = \coker \phi$ is annihilated by $a$, $b$, or $c$
(See Corollary~\ref{cor:annT-3d-characteristic-dimension-0}):
\[
 a \begin{pmatrix} b \\ a \\ 0 \end{pmatrix} =\phi \begin{pmatrix} 1 \\ 1+a \\ 0 \\ a \end{pmatrix}, \quad
 b \begin{pmatrix} b \\ a \\ 0 \end{pmatrix} =\phi \begin{pmatrix} 1 \\ 1+b \\ 1 \\ 1 \end{pmatrix}, \quad
 c \begin{pmatrix} b \\ a \\ 0 \end{pmatrix} =\phi \begin{pmatrix} 0 \\ 0 \\ 1 \\ 1 \end{pmatrix}.
\]
Therefore, $T$ is isomorphic to $\coker \partial_1 \cong \FF_2$ of Eq.~\eqref{eq:resolution-F}.
The arguments $h_x,h_y,h_z$ of $\phi$ can be thought of as \emph{hopping} operators for the charge.
According to \cite{LevinWen2003Fermions}, one can check that the charge is actually a fermion
from the commutation values among, for example, $h_x,h_y,\bar y h_y$.


Consider a short exact sequence
\[
 0 \to T \to C \to N \to 0 .
\]
The corresponding sequence for 3D toric code splits,
i.e., $C \cong T \oplus N$, while this does not.
It implies that this model is not the same as the 3D toric code.

Now we can compute the ground-state degeneracy, or $\dim_{\FF_2} K(L)$.
Tensoring the boundary condition
\[
B = R/\bb_L = R/(x^L-1,y^L-1,z^L-1)
\]
to the short exact sequence,
we have a long exact sequence
\[
 \cdots \to
\Tor_1(T,B) \xrightarrow{\delta'} \Tor_1( C, B) \xrightarrow{\delta} \Tor_1(N,B) 
\to T \otimes B \to C \otimes B \to N \otimes B \to 0 .
\]
Hence, $K(L) \cong \Tor_1( C,B)$ has vector space dimension $\dim_{\FF_2} \im \delta + \dim_{\FF_2} \ker \delta$.
Since the sequence is exact, $\dim_{\FF_2} \ker \delta = \dim_{\FF_2} \im \delta'$.
As we have seen in Example~\ref{eg:highD-toric-codes},
\begin{align*}
 \Tor_1(T,B) & \cong \Tor_1(R/\mm,B) \cong (\FF_2)^3 , \quad \text{and} \\
 \Tor_1(N,B) & \cong \Tor_2(R/\mm,B) \cong (\FF_2)^3 .
\end{align*}
It follows that $\dim_{\FF_2} K(L) \le \dim_{\FF_2} \Tor_1(N,B) + \dim_{\FF_2} \Tor_1(T,B) = 6$.

It is routine to verify that $\bb_4 \subseteq I_2(\phi) \subseteq \mm := (x+1,y+1,z+1)$.
Recall the decomposition $K(L) = \bigoplus_\pp K(L)_\pp$ where $\pp$ runs over all maximal ideals of $R/\bb_L$.
Due to Lemma~\ref{lem:associated-maximal-localized-homology}, this decomposition consists of only one summand $K(L)_\mm$.
When $L$ is odd, since $(\bb_L)_\mm = \mm_\mm$, we know $K(L)_\mm = K(1)_\mm$.
Since $\phi \mapsto 0$ under $a=b=c=0$, we see $\dim_{\FF_2} K(1) = 4$.
The logical operators in this case are
\[
 \begin{pmatrix}  0 \\ 0 \\ \widehat z  \\ \widehat z  \end{pmatrix}
\smile
 \widehat x \cdot \widehat{y \bar z}
 \begin{pmatrix}  0 \\ 1 \\ 0 \\ 0 \end{pmatrix}
%
 \quad ; \quad
%
 \begin{pmatrix}  \widehat x \\ \widehat x \\ 0 \\ 0 \end{pmatrix}
\smile
 \widehat z \cdot \widehat{x y} 
 \begin{pmatrix}  1 \\ 1 \\ 1 \\ 0 \end{pmatrix}
\]
where $\widehat \mu = \sum_{n=0}^{L-1} \mu^n$ so $\mu \cdot \widehat \mu = \widehat \mu$,
and symplectic pairs are tied.
The left elements are string-like, and the right surface-like.
% Note that the string and surface are not orthogonal.

When $L$ is even, the following are $\FF_2$-independent elements of $K(L)$.
As there are 6 in total, the largest possible number, we conclude that $K(L)$ is 6-dimensional,
i.e., the number of encoded qubits is 3 when linear dimensions are even.
\[
 \begin{pmatrix}  0 \\ 0 \\ \widehat z \\ \widehat z \end{pmatrix}
\smile
 \widehat x ' \widehat y '
 \begin{pmatrix}  1+y \\ x+xy \\ 0 \\ 1+x+y+xy \end{pmatrix}
%
 ;
%
 \begin{pmatrix}  \widehat x \\ \widehat x \\ 0 \\ 0 \end{pmatrix}
\smile
 \widehat y ' \widehat z '
 \begin{pmatrix}  1+z \\ 1+z \\ 1+z \\ y+yz \end{pmatrix}
%
  ;
%
 \begin{pmatrix}  \widehat y \\ \widehat y \\ \widehat y \\ \widehat y \end{pmatrix}
\smile 
 \widehat z ' \widehat x '
 \begin{pmatrix}  0 \\ 1+x+z+xz \\ 1+x \\ 1+x \end{pmatrix}
\]
where $\widehat \mu ' = \sum_{i=0}^{L/2-1} \mu^{2i}$ so $(1+\mu)\widehat \mu ' = \widehat \mu$.
The pairs are transformed cyclically by $x \mapsto y \mapsto z \mapsto x$
together with the symplectic transformation $\omega$ of Eq.~\eqref{eq:symplectic-Levin-Wen}.
%
\hfill $\Diamond$
\end{example}


\section{Discussion}
\label{sec:lowD-discussion}

There are many natural questions left unanswered.
Perhaps, it would be the most interesting to answer how much 
the associated ideal $I(\sigma)$
determines about the Hamiltonian.
Note that the very algebraic set defined by the associated ideal
is not invariant under coarse-graining.
For instance, in the characteristic dimension zero case,
the algebraic set can be a several points in the affine space,
but becomes a single point under a suitable coarse-graining.
However, the geometry of the algebraic set seems to be crucial
to prove, for example, ``no-strings rule.''
See Chapter~\ref{chap:cubic-code}.


It is interesting on its own to prove or disprove
that the elementary symplectic transformations 
generate the whole symplectic transformation group.
In the zero-dimensional case where $\Lambda$ is the trivial group,
it is true as we have already seen in 
Proposition~\ref{prop:elem-sym-transformation-generates-whole-BlockCodeCase}.
The one-dimensional case is also true 
because it is implied by the computation in the proof of Theorem~\ref{thm:1D-classification}.
A classical problem answered affirmatively by Suslin~\cite{Suslin1977Stability}
is that any sufficiently large invertible matrix over a polynomial ring
is a finite product of elementary matrices such as row operations and scalar multiplications.
Later, an algorithmic proof is given by Park and Woodburn~\cite{ParkWoodburn1995Algorithmic}.
A similar problem under a confusingly similar name `symplectic group' over polynomial rings
is solved by Grunewald~{\em et al.}~\cite{GrunewaldMennickeVaserstein1991SymplecticGroup},
who defined the `symplectic group' as
$ \left\{ S \in \mathrm{Mat}\left(n, \FF[x_1,\ldots,x_n]\right) ~\middle| ~ S^T \lambda S = \lambda \right\}$
where $T$ is the transpose.
Kopeyko~\cite{Kopeyko1999Symplectic} generalized it to include Laurent polynomials,
but still the `symplectic group' is different from ours
since the antipode map is absent from the definition
\[
 \left\{ S \in \mathrm{Mat}\left(n, \FF[x_1^{\pm 1},\ldots,x_n^{\pm 1}]\right) ~\middle| ~ S^T \lambda S = \lambda \right\} .
\]

The characteristic dimension is not proven to be invariant under coarse-graining.
It suffices to have an upper bound on $\dim_{\FF_2} K(L)$ in Corollary~\ref{cor:upper-bound-k}.
without the condition that $2 \nmid L$.
A closely related object is the Hilbert function.
Given a graded module $M = \bigoplus_{s=0}^{\infty} M_s$ 
over a polynomial ring with coefficients in a field $\FF$,
the Hilbert function $f_M$ is a numerical function defined by $f_M(s) = \dim_\FF M_s$.
Since $K(L) = \dim_{\FF_2} \Tor_1(\coker \epsilon, R / \bb_L)$,
the Hilbert function might be useful if we could make $\coker \epsilon$ graded.
A technical difficulty would be that the ideal $\bb_L=(x_1^L-1,\ldots,x_D^L-1)$ 
is not a power of $\mm=(x_1 - 1,\ldots,x_D-1)$.
See \cite[Chapter~12]{Eisenbud} and \cite{IyengarPuthenpurakal2007HilbertSamuelFunctions}.

Lastly, an important problem is to give a criterion to a module $M$
that can be realized as $\coker \sigma^\dagger$
for an exact code Hamiltonian described by $\sigma$.
In the two-dimensional case, we know the answer ---
$\ann M = (x-1,y-1)$ by Theorem~\ref{thm:structure-2d-ann-coker-epsilon}.
%(include?) The conjecture in Section~\ref{sec:extracting-toric-codes} speculates the uniqueness of resolutions of $M$.